\documentclass{ropas-poster}

\usepackage{booktabs}

\title{Example for \texttt{ropas-poster} class}
\subtitle{Basic example}
\author{SHIN Chaehwan}
\institute{Programming Research Laboratory (ROPAS), Seoul National University}
\logos{../.images/logo-snu-eng.png,../.images/logo-ropas.png}

\begin{document}
  \begin{postergrid}{1}{2}
    \begin{posterblock*}{1}{1}
      \begin{postergrid}{2}{1}
        \begin{posterblock}{1}{1}
          \blocktitle{1. Mathematical expressions}

          If \begin{math}
            {x}\mathrel{\leq}{y}
          \end{math} and \begin{math}
            {y}\mathrel{\leq}{z}
          \end{math}, then \begin{math}
            {x}\mathrel{\leq}{z}
          \end{math} for all \begin{math}
            {x,y,z}\mathrel{\in}{X}
          \end{math}.
          That is, \begin{displaymath}
            {\forall}{{x,y,z}\mathrel{\in}{X}}.\,{{{{x}\mathrel{\leq}{y}}\mathrel{\land}{{y}\mathrel{\leq}{z}}}\mathrel{\rightarrow}{{x}\mathrel{\leq}{z}}}.
          \end{displaymath}
        \end{posterblock}

        \begin{posterblock}{2}{1}
          \blocktitle{2. Lists}

          Let \begin{math}
            G
          \end{math} be a set, \begin{math}
            {\bullet}\mathrel{:}{{{G}\mathbin{\times}{G}}\mathbin{\rightarrow}{G}}
          \end{math}, \begin{math}
            {{\left(\cdot\right)}^{-1}}\mathrel{:}{{G}\mathbin{\rightarrow}{G}}
          \end{math}, and \begin{math}
            {e}\mathrel{\in}{G}
          \end{math}.
          If \begin{itemize}
            \item
            \begin{math}
              \bullet
            \end{math} is associative,
            \item
            \begin{math}
              {{e}\mathbin{\bullet}{x}}\mathrel{=}{x}
            \end{math} and \begin{math}
              {{x}\mathbin{\bullet}{e}}\mathrel{=}{x}
            \end{math} for all \begin{math}
              {x}\mathrel{\in}{G}
            \end{math}, and
            \item
            \begin{math}
              {{{x}^{-1}}\mathbin{\bullet}{x}}\mathrel{=}{e}
            \end{math} and \begin{math}
              {{x}\mathbin{\bullet}{{x}^{-1}}}\mathrel{=}{e}
            \end{math} for all \begin{math}
              {x}\mathrel{\in}{G}
            \end{math},
          \end{itemize} then we call \begin{math}
            \left({G},{\bullet},{{\left(\cdot\right)}^{-1}},{e}\right)
          \end{math} a \textbf{group}.

          Let \begin{math}
            G
          \end{math} be a set and \begin{math}
            {\bullet}\mathrel{:}{{{G}\mathbin{\times}{G}}\mathbin{\rightarrow}{G}}
          \end{math}.
          If \begin{enumerate}
            \item
            \begin{math}
              \bullet
            \end{math} is associative, and
            \item
            there exists \begin{math}
              {e}\mathrel{\in}{G}
            \end{math} such that \begin{enumerate}
              \item
              \begin{math}
                {{e}\mathbin{\bullet}{x}}\mathrel{=}{x}
              \end{math}, and
              \item
              there exists \begin{math}
                {y}\mathrel{\in}{G}
              \end{math} such that \begin{math}
                {{y}\mathbin{\bullet}{x}}\mathrel{=}{e}
              \end{math} for all \begin{math}
                {x}\mathrel{\in}{G}
              \end{math},
            \end{enumerate} then there uniquely exists a group with \begin{math}
              G
            \end{math} and \begin{math}
              \bullet
            \end{math}.
          \end{enumerate}
        \end{posterblock}
      \end{postergrid}
    \end{posterblock*}

    \begin{posterblock*}{1}{2}
      \begin{postergrid}{3}{1}[row-ratio={1,min,min}]
        \begin{posterblock}{1}{1}
          \blocktitle{3. Tables}

          \begin{table}
            \caption{Small table}
            \begin{tabular}{ccc}
              \toprule
              Row 1 & Row 2 & Row 3 \\
              \midrule
              Item 1 & Item 2 & Item 3 \\
              Item 4 & Item 5 & Item 6 \\
              \bottomrule
            \end{tabular}
          \end{table}

          Lorem ipsum dolor sit amet, consectetur adipiscing elit, sed do eiusmod tempor incididunt ut labore et dolore magna aliqua.
          Ut enim ad minim veniam, quis nostrud exercitation ullamco laboris nisi ut aliquip ex ea commodo consequat.
          Duis aute irure dolor in reprehenderit in voluptate velit esse cillum dolore eu fugiat nulla pariatur.
          Excepteur sint occaecat cupidatat non proident, sunt in culpa qui officia deserunt mollit anim id est laborum.

          \begin{table}
            \caption{Big table}
            \begin{tabular}{crcl}
              \toprule
              Symbol & English & & German \\
              \midrule
              \begin{math}
                0
              \end{math} & zero & \begin{math}
                \leftrightarrow
              \end{math} & Null \\
              \begin{math}
                1
              \end{math} & one & \begin{math}
                \leftrightarrow
              \end{math} & Eins \\
              \begin{math}
                2
              \end{math} & two & \begin{math}
                \leftrightarrow
              \end{math} & Zwei \\
              \begin{math}
                3
              \end{math} & three & \begin{math}
                \leftrightarrow
              \end{math} & Drei \\
              \begin{math}
                4
              \end{math} & four & \begin{math}
                \leftrightarrow
              \end{math} & Vier \\
              \begin{math}
                5
              \end{math} & five & \begin{math}
                \leftrightarrow
              \end{math} & F\"unf \\
              \begin{math}
                6
              \end{math} & six & \begin{math}
                \leftrightarrow
              \end{math} & Sechs \\
              \begin{math}
                7
              \end{math} & seven & \begin{math}
                \leftrightarrow
              \end{math} & Sieben \\
              \begin{math}
                8
              \end{math} & eight & \begin{math}
                \leftrightarrow
              \end{math} & Acht \\
              \begin{math}
                9
              \end{math} & nine & \begin{math}
                \leftrightarrow
              \end{math} & Neun \\
              \bottomrule
            \end{tabular}
          \end{table}
        \end{posterblock}

        \begin{posterblock}{2}{1}
          \blocktitle{4. Minimum cell}

          Lorem ipsum dolor sit amet, consectetur adipiscing elit, sed do eiusmod tempor incididunt ut labore et dolore magna aliqua.
          Ut enim ad minim veniam, quis nostrud exercitation ullamco laboris nisi ut aliquip ex ea commodo consequat.
          Duis aute irure dolor in reprehenderit in voluptate velit esse cillum dolore eu fugiat nulla pariatur.
          Excepteur sint occaecat cupidatat non proident, sunt in culpa qui officia deserunt mollit anim id est laborum.
        \end{posterblock}

        \begin{posterblock}{3}{1}
          \blocktitle{5. Thank you!}
        \end{posterblock}
      \end{postergrid}
    \end{posterblock*}
  \end{postergrid}
\end{document}
